\section{Introduction}

The advancements in Large Language Models (LLMs), coupled with the wide availability of LLM-based conversational agents, have democratized the power of AI.
Given the open-ended nature of LLMs, the purposes, contexts, and consequences of using them become more directly determined by the end-users.
\revision{As a result, users are playing an increasingly important role in conforming to responsible AI principles such as fairness, transparency and privacy~\cite{diaz2023connecting, dignum2023responsible}, and ensuring AI is not used in a way that causes harms to humanity.}

Transparency, as one of the responsible AI principles, has been significantly impacted by LLMs~\cite{liao2023ai}.
\revision{Prior literature on AI transparency has focused on transparency via explainability techniques, regulatory or legal frameworks, or data-usage disclosures, attributing the responsibility of providing transparency to AI practitioners, researchers, or other technical experts~\cite{corbett2023interrogating}.
However, the democratization of AI, especially LLMs, has extended this responsibility to \textit{AI users}, raising transparency requirements regarding AI usage.}
Due to the potential issues in LLM-generated content, including bias~\cite{abid2021persistent, isaza2023fairy}, hallucination~\cite{bender2021dangers, dahl2024large, yao-etal-2024-samples}, invasion of intellectual property~\cite{karamolegkou2023copyright}, and threats to human autonomy~\cite{parasuraman2010complacency}, it has become crucial for everyday users of LLM-based services to honestly and accurately report when and how LLM-generated content is used in real life.
This is not only important for cautioning against and preventing potential negative consequences caused by LLMs, but also serves as a prerequisite for studying and mitigating these negative impacts.
Disclosure of the use of LLMs has become a basic requirement for \textit{AI users} in various domains to adhere to, such as academic research~\cite{ACM2023AuthorshipPolicy, Elsevier2024AIWritingPolicy}, education~\cite{Harvard2024ChatGPTGuidelines, UMNLib2024LibGuide, UBCLibrary2024CitingGenAI, MSU2024GenAIGuidance}, freelance jobs~\cite{Upwork2024ChatGPTStance}, and online publishing~\cite{Medium2024DistStandards}.
For example, academic publishers like ACM~\cite{ACM2023AuthorshipPolicy} and Elsevier~\cite{Elsevier2024AIWritingPolicy} permit the use of generative AI such as ChatGPT but require full disclosure on how AI has been used for creating content in the publications, given that a lack of transparency in such instances might result in overlooked issues concerning the scientific rigor of the published work~\cite{kaebnick2023editors}.


On the other hand, a recent study~\cite{zhang2024s} on LLM-based conversational agents discovered a novel privacy concern focusing on the users' apprehension that \textit{the action of using LLM-based conversational agents} in a certain context may be discovered.
For example, the paper documented a case in which a non-native English speaker used ChatGPT for revising emails and was worried that her professors might ``change their attitudes to her'' once discovering her reliance on AI and hoped ``they will never know I used AI to do that (write emails).''
This fear could ultimately result in a tendency to hide the use of AI and keep it a secret from others, suggesting a general problem concerning the tension between privacy needs of the individual and transparency needs for the greater good.

However, aside from a few examples, \textbf{there is still a lack of systematic understanding of the landscape, mechanisms, and implications of secret use of LLMs}. Without such understanding, it is infeasible to gauge its potential impact, identify mitigation strategies when privacy and transparency are at odds with each other, and design effective interventions to encourage transparent disclosure of LLM use.
To bridge this gap, we aim to thoroughly study the phenomenon of the secret use of LLMs.
We seek to answer a fundamental question: \textit{When and why do people hide their use of LLMs\footnote{In this paper, we define the scope of ``use of LLMs'' to cover a wide range of activities that involve using content generated by LLMs in real-world tasks. The interaction paradigm might involve widely-used LLM-based conversational agents such as ChatGPT, as well as those used by tech-savvy users, such as open-source LLMs like Llama. We have both types of users in our sample.}?} 

We took a mixed-methods approach to answer this research question, using two sequential studies. 
Given the limited examples found in prior research, we first carried out an online survey with \surveyallusernumber{} LLM users to gain a holistic understanding of the landscape of the phenomenon. 
Through qualitative analysis of \surveycasenumber{} real-world cases of secret LLM usage reported by the users, we summarized the contexts in which secret use of LLMs occurred and developed a typology of causes of users' secret use of LLMs.
We observed that secret LLM usage happened in critical domains such as academic writing and work tasks, and the causes can mainly be divided into two categories -- how users view their usage of LLMs (internal judgement), and how users perceive others view it (perceived external judgement).


Informed by the survey results, we conducted a controlled experiment study with 300 LLM users to quantitatively measure the variance of secret LLM usage among different types of tasks, and investigate the mechanism behind the secret LLM usage behavior.
We focused on understanding the factors influencing the \textit{intention to perform the secret use behavior} at two levels: \revision{passive non-disclosure, where users do not actively disclose their use of LLMs but are honest if asked, and active concealment, where users deliberately hide their use of LLMs even when being directly inquired about it.}
Compared to the baseline task, ``General Information Search,'' we found that the five experimental tasks - ``Creative Writing,'' ``Academic Writing,'' ``School Works,'' ``Work Tasks,'' and ``Social Connection'' - all resulted in significantly higher intentions of both passive non-disclosure (40\% vs. 72\%) and active concealment (16\% vs. 54\%).
However, we did not observe a significant impact due to individual differences, which suggests that the secret use behavior generally applies to people with different demographic backgrounds and personality traits.
Our mediation analysis further reveals that task types primarily affect the intention to use secretive behavior by influencing perceived external judgment.
This provides important insights for future work on designing interventions to encourage more transparent disclosure of LLM/AI use.

With the findings of the two studies, we further discuss the societal and individual-level implications of the secret use of LLM behaviors.
This includes the different priorities between privacy and transparency in various contexts, as well as the potential emotional stress caused to the user by the secretive behavior.
We conclude with a set of suggestions for stakeholders to build community norms and improve community cultures, in order to mitigate the barriers to transparent disclosure caused by perceived external judgement.
We also identify future research directions to further explore this topic.

